\section{Conclusion}

We studied directional scale inconsistency in chart question answering. The main empirical finding is not that one adapter universally fixes scale errors, but that correction is regime-dependent. In answer-tuned Qwen2.5-VL-7B, large-positive scale-contrastive adaptation gives a small held-out ChartQA gain and recovers nearly half of strict test-set scale errors. ScaleQA-Natural explains the boundary: self-check is stronger than a global adapter, small-positive policies help different regimes, and aggressive selectors can hurt by switching too often.

The strongest evidence is the gap between candidate availability and candidate choice. Structured verifiers, likelihood scoring, pairwise comparison, and top-8 pools often contain a correct scaled answer, yet deployable selectors fail to choose it. Our best no-leak selector improves over self-check but remains far below policy and candidate oracles. Cross-backbone and external checks show that scale-like failures transfer more readily than a correction recipe.

Overall, Chart QA systems should evaluate scale reliability by regime, with fix--break accounting and candidate-selection analysis alongside aggregate accuracy.
